\documentclass[11pt]{article}

\usepackage[margin=1in]{geometry}
\usepackage[utf8]{inputenc}
\usepackage[T1]{fontenc}
\usepackage{amsmath}
\usepackage{booktabs}
\usepackage{graphicx}
\usepackage{hyperref}
\usepackage{natbib}
\usepackage{xcolor}

\title{Identification-First Causal Inference for Pacific Northwest Extreme Precipitation}
\author{Pradyumna}
\date{UCSD DSC 245 Course Project}

\begin{document}
\maketitle

\begin{abstract}
This project studies whether warm regional sea-surface temperature (SST)
anomalies increase cool-season extreme-precipitation risk in the Pacific
Northwest under an explicit causal identification graph. Using ERA5 monthly
reanalysis from 1979--2023, the pipeline combines PCMCI+ and VARLiNGAM
discovery, a documented projection to a lag-0 acyclic panel DAG, a domain
identification backbone, IPW and doubly robust ACE estimation, falsification
checks, temporal holdout validation, logistic baselines, and sensitivity
analyses. The headline ERA5 result is cautiously positive for the Pacific
Northwest ONDJFM estimand: full-sample IPW estimates a risk difference of 0.146
and DR estimates 0.160; temporal holdout estimates are also positive. The claim
does not generalize to California, the Intermountain West, pooled western-US
stacking, or the NDJF-only seasonal variant. The main contribution is an
auditable causal workflow for narrow climate claims, not a new estimator.
\end{abstract}

\section{Question and Decision Context}

The central question is narrow: under an explicit identification graph, does
warm regional SST increase ONDJFM 90th-percentile precipitation-extreme risk in
the Pacific Northwest?

The practical decision context is seasonal risk screening. Pacific Northwest
water managers, flood-risk planners, reservoir operators, and emergency
preparedness offices can use a validated risk-difference estimate as one input
to scenario stress tests and monitoring priorities. This is not an operational
forecast and not attribution for individual storms.

\section{Data and Panel Construction}

The project uses ERA5 monthly reanalysis \citep{hersbach2020era5} from
1979-01 through 2023-12. The western-US domain is aggregated into three regional
panels: Pacific Northwest, California, and Intermountain West. The Pacific
Northwest box is 45--50 N and 125--115 W. Regional SST means the monthly SST
anomaly averaged over that same region box.

\begin{figure}[h]
\centering
\includegraphics[width=0.75\linewidth]{figures/fig1_domain_map.pdf}
\caption{Western-US analysis domain and regional masks. The primary claim is
Pacific Northwest ONDJFM.}
\end{figure}

The treatment is \texttt{sst\_warm}, equal to one when regional SST anomaly is at
least 0.5 K. The outcome is \texttt{tp\_extreme}, equal to one when regional
monthly precipitation is at or above the region-specific 90th percentile. The
primary season is ONDJFM: October, November, December, January, February, and
March.

\section{Discovery-to-Identification Workflow}

PCMCI+ \citep{runge2019detecting,runge2020discovering} and VARLiNGAM
\citep{shimizu2006linear,hyvarinen2010estimation} are used as discovery tools,
not as proof of the true climate graph. Their outputs are converted into a
monthly panel graph through a documented projection:

\begin{enumerate}
\item Drop self-edges because they do not define cross-variable panel
adjustment.
\item Retain lag-0 links only for the panel DAG used by DoWhy and ACE
estimation.
\item Sort candidate lag-0 edges by discovery strength.
\item Greedily add edges while preserving acyclicity.
\item Merge the resulting structural panel DAG with a domain identification
backbone.
\end{enumerate}

The domain backbone is:
\[
\texttt{nino34}\rightarrow\texttt{z500},\quad
\texttt{nino34}\rightarrow\texttt{sst},\quad
\texttt{z500}\rightarrow\texttt{tp},\quad
\texttt{sst}\rightarrow\texttt{tp},\quad
\texttt{tp}\rightarrow\texttt{tp\_extreme}.
\]

This is the key modeling choice. It makes the assumptions visible, but it also
means the causal interpretation depends on the backbone and the projection rule.

\section{Identification and Adjustment}

Under the merged identification graph, the primary backdoor adjustment set for
\texttt{sst} to \texttt{tp\_extreme} in the Pacific Northwest is
\{\texttt{z500}\}. This is intentionally minimal. Adding variables is not
automatically better: \texttt{nino34}, \texttt{t2m}, and \texttt{swvl1} can play
different roles under different graphs, including ancestors, mediators,
descendants, or precision variables. The project therefore treats
\{\texttt{z500}\} as the primary graph-identified set and uses larger observed
sets as sensitivity checks.

\section{Primary Results}

\begin{figure}[h]
\centering
\includegraphics[width=\linewidth]{figures/fig7_primary_holdout_summary.pdf}
\caption{Main risk-difference estimates. The Pacific Northwest signal is
positive in the full sample and temporal holdout; California, the Intermountain
West, and pooled western-US estimates are weak or null.}
\end{figure}

\begin{table}[h]
\centering
\caption{Primary Pacific Northwest ONDJFM estimates.}
\begin{tabular}{lrr}
\toprule
Estimator & Risk difference & 95\% CI \\
\midrule
IPW & 0.146 & [0.026, 0.237] \\
DR & 0.160 & [0.040, 0.245] \\
\bottomrule
\end{tabular}
\end{table}

Overlap diagnostics are acceptable for the primary IPW model: ESS/n = 0.701,
propensity range [0.116, 0.408], and 99th-percentile weight = 6.217.

\section{Temporal Holdout}

The holdout fits nuisance models on 1979--1999 and evaluates on 2000--2023.

\begin{table}[h]
\centering
\caption{Temporal holdout estimates.}
\begin{tabular}{lrr}
\toprule
Estimator & Risk difference & 95\% CI \\
\midrule
IPW holdout & 0.180 & [0.016, 0.361] \\
DR holdout & 0.066 & [0.045, 0.085] \\
\bottomrule
\end{tabular}
\end{table}

The DR interval is much tighter than IPW. The audit shows the IPW interval is
8.7 times wider. This is not a free robustness guarantee: both holdout
estimators keep train-fitted nuisance models fixed during evaluation bootstrap
resampling, and the DR estimate is stabilized by its outcome-effect model.

\section{Baseline Logistic Comparison}

\begin{table}[h]
\centering
\caption{Naive and adjusted logistic baselines.}
\begin{tabular}{llrr}
\toprule
Model & Covariates & Risk difference & 95\% CI \\
\midrule
Crude logistic & \texttt{sst\_warm} & 0.094 & [-0.008, 0.205] \\
Graph-adjusted logistic & \texttt{sst\_warm + z500} & 0.147 & [0.051, 0.237] \\
Expanded logistic & \texttt{sst\_warm + z500 + nino34 + t2m + swvl1} & -0.007 & [-0.105, 0.100] \\
\bottomrule
\end{tabular}
\end{table}

The adjusted logistic result is close to IPW, which means the signal is visible
in a simple model. The value of the causal workflow is that it formalizes the
estimand, records graph assumptions, checks overlap, runs falsifications, and
performs temporal validation.

\section{Backbone and Lag Sensitivity}

Graph perturbations that add \texttt{nino34 -> tp}, add \texttt{z500 -> sst}, or
remove \texttt{nino34 -> sst} preserve a positive primary estimate. However,
broad observed adjustment is null and has poor overlap: ESS/n falls to 0.198.
This is the strongest remaining weakness in the identification story and should
be framed as graph dependence.

\begin{table}[h]
\centering
\caption{Lag/projection sensitivity.}
\begin{tabular}{lrrl}
\toprule
Variant & IPW risk difference & 95\% CI & Note \\
\midrule
Base \{\texttt{z500}\} & 0.145 & [0.036, 0.246] & Positive \\
Add \texttt{sst\_lag1} & 0.166 & [-0.056, 0.441] & Positive but imprecise \\
Add lagged circulation/ENSO & 0.227 & [0.047, 0.449] & Positive \\
Add \texttt{tp\_lag1} & 0.156 & [0.037, 0.282] & Positive stress check \\
\bottomrule
\end{tabular}
\end{table}

The result is not solely an artifact of dropping lagged variables, although some
lagged specifications are less stable.

\section{Negative Results and Seasonal Provenance}

California, the Intermountain West, and pooled western-US stacking are weak or
null. The NDJF-only seasonal variant also fails sign stability. ONDJFM is the
frozen primary season for the current analysis release and is climatologically
motivated as a broad Pacific Northwest cool-season window, but the repository
does not prove fully prospective seasonal pre-registration. Temporal holdout
validates the current frozen claim but does not remove all seasonal-selection
concern.

\section{Limitations}

The causal interpretation depends on the identification graph, especially the
domain backbone and projected lag-0 graph. ERA5 is a reanalysis product, not
direct observational truth. Monthly aggregation smooths submonthly storm
dynamics. Atmospheric rivers are likely mechanisms behind many Pacific Northwest
precipitation extremes \citep{dettiger2011climate}, but this monthly panel does
not include atmospheric-river covariates and should not be read as storm-scale
AR causality. External validation in MERRA-2 or JRA-55 remains future work.

\section{Final Claim}

Under a frozen identification graph combining a domain backbone with a lag-0
acyclic projection of causal-discovery output, warm regional SST anomalies
increase ONDJFM 90th-percentile precipitation-extreme risk in the Pacific
Northwest in ERA5. The effect is positive under IPW and DR and strictly
replicates in a post-2000 temporal holdout, but it is geographically,
seasonally, and graph-assumption specific.

\bibliographystyle{plainnat}
\bibliography{references}

\end{document}
